主码 ADC 使用 FP32 $q'$ 与存储的重建 \cite{pq,jhq}：
\begin{equation}
 T_m[z]=\|q'_m-\mathcal C_m[z]\|_2^2,\qquad
 \primarydist(q,x)=\sum_{m=0}^{M-1}T_m[z_m(x)].
 \label{eq:5-primary-lut}
\end{equation}
因为每个子码本具有 Cartesian 结构，256 项表可以缩小。

对 $K=256$、$D_s=8$，解析码本给出
\begin{equation}
 \begin{aligned}
 a&=\sigma\Phi^{-1}(3/4)\simeq0.67449\sigma,\qquad
 \mathcal C_m=\{-a,+a\}^{8},\\
 \primarydist(q,x)&=\|q'\|_2^2+da^2
       -2a\sum_{j=0}^{d-1}s_j(x)q'_j.
 \end{aligned}
 \label{eq:5-sign-codebook}
\end{equation}
其中 $\Phi$ 是标准正态 CDF；一个 byte 存八个符号，高位在前。因此
\[
 a>0\quad\Longrightarrow\quad
 \arg\min_x\primarydist(q,x)=\arg\max_x\sum_j s_j(x)q'_j.
\]
四路／八路拆分保持相同的实数算术分数。学习式替代码本需要兼容性检查。

\subsection{不改变量化器的分解}

写成 $z=16h+l$，其中 $h,l\in\{0,\ldots,15\}$、$\mathcal C_m[z]=(A_m[h],B_m[l])$。于是
\begin{equation}
 T_m^{\rm hi}[h]=\|q'_{m,\rm hi}-A_m[h]\|_2^2,\qquad
 T_m^{\rm lo}[l]=\|q'_{m,\rm lo}-B_m[l]\|_2^2.
 \label{eq:5-lut-halves}
\end{equation}
由坐标分离可得
\begin{equation}
 T_m[z]=T_m^{\rm hi}[z\!\gg\!4]
        +T_m^{\rm lo}[z\mathbin{\&}15].
 \label{eq:5-lut-factorisation}
\end{equation}
存储降为 $16+16=32$ 项；编码、候选及实数算术分数保持不变（\cref{fig:5-cartesian-lut}）。

\par\medskip\noindent\begin{minipage}{\linewidth}
\captionsetup{hypcap=false}
\centering
% Schematic indexing paths; exact identity is conditional on Cartesian coordinates.
\begin{tikzpicture}[x=1cm,y=1cm,font=\footnotesize\linespread{1}\selectfont,
 box/.style={draw=ThesisNavy!60,rounded corners=2pt,fill=ThesisNavy!4,inner sep=4pt,align=center},
 arr/.style={-{Stealth[length=4pt]},semithick,ThesisNavy!80}]
\node[font=\small] at (7.3,5.40)
 {Cartesian premise: $\mathcal C_m[z]=(A_m[h],B_m[l])$};
\node[box,fill=ThesisNavy!12] (byte) at (7.3,4.50) {$z=(hhhh\ llll)_2=16h+l$};
\node[anchor=west,font=\bfseries] at (.15,4.95) {Full table};
\node[anchor=east,font=\bfseries] at (14.7,4.95) {Two small tables};
\node[box] (fullindex) at (2.8,3.35) {index $z\in[0,255]$};
\node[box] (hiindex) at (9.25,3.35) {$h=z\gg4$\\$hhhh$: four signs\\$2^4=16$ patterns};
\node[box] (loindex) at (12.65,3.35) {$l=z\mathbin{\&}15$\\$llll$: four signs\\$2^4=16$ patterns};
\draw[arr] (byte.west)-|(fullindex.north);
\draw[arr] (byte.east)-|(loindex.north);
\draw[arr] (byte.south)-|(hiindex.north);

\draw[ThesisNavy!55] (1.05,1.40) rectangle (4.55,2.65);
\node at (2.8,2.42) {$T_m[0],\ldots,T_m[255]$};
\draw[ThesisNavy!35] (1.05,2.17)--(4.55,2.17);
\node[fill=ThesisNavy!12,minimum width=3.4cm,minimum height=.42cm]
 (fullread) at (2.8,1.87) {$T_m[z]$};
\draw[arr] (fullindex.south)--(2.8,2.65);
\node[align=center] at (2.8,.88) {256 entries\quad 1 read\\1,024 bytes (FP32)};

\foreach \xx/\name/\idx in {9.25/hi/h,12.65/lo/l} {
 \draw[ThesisNavy!55] (\xx-1.45,1.40) rectangle (\xx+1.45,2.65);
 \node at (\xx,2.42) {$T_m^{\rm \name}[0..15]$};
 \draw[ThesisNavy!35] (\xx-1.45,2.17)--(\xx+1.45,2.17);
 \node[fill=ThesisNavy!12,minimum width=2.8cm,minimum height=.42cm]
  (read\name) at (\xx,1.87) {$T_m^{\rm \name}[\idx]$};
}
\draw[arr] (hiindex.south)--(9.25,2.65);
\draw[arr] (loindex.south)--(12.65,2.65);
\node[draw=ThesisNavy,circle,inner sep=2pt] (sum) at (10.95,.88) {$+$};
\draw[arr] (9.25,1.40)|-(sum.west);
\draw[arr] (12.65,1.40)|-(sum.east);
\node[align=center] at (10.95,.10) {16+16 entries\quad 2 reads + 1 sum\\128 bytes (FP32): $1/8$ of full-table storage};
\draw[dashed,ThesisNavy!40] (6.25,.1)--(6.25,3.45);
\node[align=center] at (7.3,-.86)
 {Same byte codes and routed candidates; equal scores in real arithmetic.\\
 FP32 evaluation can differ in rounding and ties (T4, \cref{sec:7-3}).};
\end{tikzpicture}

\captionof{figure}[Cartesian 主码 LUT 分解]{Cartesian 主码 LUT 分解示意。一个 byte 选择 256 项表中的一项，或用高／低 nibble 分别选择两个 16 项表中的一项。坐标分离在实数运算下给出加法恒等式；FP32 舍入可能改变 top-$c$ 和最终 ID。T4 中各表形式的最大相对分数差为 $4.64\times10^{-7}$；Vogue 有 2--4/256 个查询的最终 ID 不同，arXiv／OpenAI-3072 没有差异（\cref{sec:7-3}）。读取和求和描述的是查表表达式。}
\label{fig:5-cartesian-lut}
\end{minipage}\par\medskip

构造器切分码字 $16h,l$；\texttt{verify\_split\_lut} 检查 nibble／切片兼容性。任意学习得到的 256 质心 PQ 子码本不一定满足该条件：将编码标签分成 nibbles，并不使质心坐标自动可分。

\cref{sec:7-3} 的同候选 FP32 审计测量不同表形式的数值与排序差异，包括主码保留集合相同后出现的变化。

\subsection{分组、构造工作量与资源成本}

表分组消融将拆分推广到 $\glut\in\{1,2,4,8\}$ 个连续坐标组，其 bit 字段与 Cartesian 数字对应。对八位及 $D_s=8$，每组有 $2^{8/\glut}$ 项；每查询的总项数和 FP32 存储为
\begin{equation}
 n_{\rm LUT}(\glut)=\glut 2^{8/\glut},\qquad S_{\rm LUT}(\glut)=4Mn_{\rm LUT}(\glut)\ \text{bytes}.
 \label{eq:5-lut-cost}
\end{equation}
独立构造表项时，每子空间计算 $D_s2^{8/\glut}$ 个坐标贡献。因此，相较逐项独立构造整表，两组将表存储减为八分之一、坐标工作量减为十六分之一。这是表级成本比，不是应用加速预测。

\par\medskip\noindent\begin{minipage}{\linewidth}
\captionsetup{hypcap=false}
\centering
\captionof{table}[Cartesian 查找表的权衡]{八位、八坐标主码子空间的 Cartesian 表权衡。合并加法是连接查得值所需的代数最小值，不含跨子空间累加；不是实测指令计数。}
\label{tab:5-lut-resources}
\small
\begin{tabular}{@{}rrrrrr@{}}
\toprule
$\glut$ & 项数 & FP32 字节 & 查表数 & 合并加法 & 构造项数 \\
\midrule
1 & 256 & 1,024 & 1 & 0 & 2,048 \\
2 & 32 & 128 & 2 & 1 & 128 \\
4 & 16 & 64 & 4 & 3 & 32 \\
8 & 16 & 64 & 8 & 7 & 16 \\
\bottomrule
\end{tabular}
\end{minipage}\par\medskip

表先在全局工作区构建；若完整共享内存占用符合查询到的每 block 可选上限 $S_{\max}$，则协作暂存到 shared memory。缓冲容量为 $C_{\rm buf}$ 时，普通／分组扫描占用为
\begin{equation}
 S_0=8C_{\rm buf}+8+S_{\rm LUT}(\glut),\qquad
 S_{\rm grp}=8C_{\rm buf}+24+S_{\rm LUT}(\glut)
 \quad\text{bytes}.
 \label{eq:5-shared-footprint}
\end{equation}
对已评测路径（不含可选融合精排／提前退出临时区），
\[
 \text{shared lookup}\iff S_{\rm path}\leq S_{\max}.
\]

表扫描实验中，$c=1{,}000$，block 大小为 512 或 1,024，$C_{\rm buf}=2{,}048$。整表占 96--384\,KiB，两组为 12--48\,KiB。NVIDIA 对计算能力 12.0 规定每 block 上限 99\,KiB \cite{blackwell_tuning}。按此上限，即使 $M=96$，整表占用仍超容量（query-centric 路径为 $114{,}696$ bytes），而最大两组占用在分组路径中为 $65{,}560$ bytes。因此，这些配置下的源码放置规则将整表放在 global memory，将两组表放在 shared memory。这是依据启动规则和公开上限推断的放置；测量记录未保存设备属性查询。

四组和八组每子空间均存 16 项，但查表和合并更多。位提取、依赖、暂存、同步和驻留共同决定权衡。每个扫描组可能暂存一份副本。

$16+16$ nibble LUT 沿用小表先例 \cite{fastscan,quickeradc,ivfrabitq}。本文通过 GPU 保留／精排将其与 JHQ 两级分数集成；实测权衡见 \cref{sec:7-3}。
